\subsection{2026-09-23 --- Hyperparameters, and what effective sample size means in two shapes}
\label{sec:tuning}

\subsubsection*{What we did}

Neither model had been tuned. The pointwise model's settings were inherited from the wide
architecture; the count model's were unchanged since Section~\ref{sec:oracle}, surviving a
3-class $\to$ 6-class head change, growth to 462 features, and the entire long-format rebuild.

Coordinate descent, one knob at a time, validation year 2021, stride 3 held constant. Not a final
tuning phase --- a 20-minute probe to find anything badly mis-set.

\subsubsection*{The pointwise model: $+0.0028$}

\begin{center}\small
\begin{tabular}{lrr}
\toprule
knob & from $\to$ to & micro-AUC \\
\midrule
baseline & --- & 0.9604 \\
\texttt{min\_child\_weight} & $1 \to 500$ & \textbf{0.9622} \\
\texttt{max\_depth} & $8 \to 10$ / $12$ & 0.9621 / 0.9620 (no gain) \\
capacity & $300{\times}0.10 \to 600{\times}0.05$ & \textbf{0.9632} \\
\texttt{colsample\_bytree} & $0.8 \to 0.5$ / $0.3$ & 0.9625 / 0.9626 (no gain) \\
\bottomrule
\end{tabular}
\end{center}

\begin{finding}
\texttt{min\_child\_weight} was the one badly mis-set knob, and it was predictable from the data
rather than from a grid. The long-format table has 5.2M rows, but
$P(\text{on}\mid\text{on}_{t-1}) = 0.995$ and each unit-week holds 1--3 contiguous blocks
(Section~\ref{sec:temporal}), so the independent information is roughly 40k--120k blocks. A leaf
minimum of 1 lets a split form from rows that are all the same block of the same unit-week ---
fitting autocorrelation. Raising it to 500 is worth $+0.0018$.
\end{finding}

Depth not helping is worth recording as a wrong prediction: the argument was that long format
exists to allow per-unit specialisation and needs depth to pay for it. It does not, because the
\texttt{cf\_*} features are already unit-specific, so the tree never has to split on unit first
to reach unit-specific structure.

\subsubsection*{The count model: nothing beats the original settings}

Screened on multi-class \textbf{log-loss}, not count MAE. Section~\ref{sec:abcd} established MAE
moves opposite to micro-AUC because it scores only the collapsed point estimate, while what is
forwarded is the whole distribution; log-loss is a proper scoring rule for that distribution.

\begin{center}\small
\begin{tabular}{lrr}
\toprule
& log-loss & count MAE \\
\midrule
\textbf{current} (400{\texttimes}0.05, d5, mcw 1, cs 0.6) & \textbf{1.0989} & 23.48h \\
\texttt{min\_child\_weight} 10 / 50 / 200 & 1.1203 / 1.1087 / 1.1517 & \\
\texttt{max\_depth} 4 / 6 / \textbf{8} & 1.1207 / 1.1214 / \textbf{1.1983} & \textbf{20.69h} \\
800{\texttimes}0.03 / 1200{\texttimes}0.02 & 1.1091 / 1.0994 & \\
\bottomrule
\end{tabular}
\end{center}

\begin{finding}
No parameter beats the settings chosen in Section~\ref{sec:oracle}. Downstream confirmation:
old params 0.9632, ``tuned'' params 0.9632 --- byte-identical, because the search returned the
starting point.
\end{finding}

\begin{finding}[The clearest demonstration of the MAE trap]
\texttt{max\_depth=8} gives the \textbf{best count MAE (20.69h) and the worst log-loss (1.1983)}.
Optimising MAE would have selected precisely the wrong model. Every MAE-guided decision from
Section~\ref{sec:oracle} onward is suspect in this direction.
\end{finding}

\begin{finding}[Same principle, opposite answers]
\texttt{min\_child\_weight} is best at \textbf{500} for the pointwise model and \textbf{1} for
the count model. Not a contradiction: count rows are whole unit-weeks and genuinely independent,
so effective $n \approx$ nominal $n$ and small leaves are legitimate. Pointwise rows are hourly
and autocorrelated at 0.995, so effective $n$ is $\sim$50$\times$ below nominal. The knob should
be set from the correlation structure of the rows, not from a default.
\end{finding}

\subsubsection*{The forwarded-feature budget is also already right}

The count model's top 120 features by gain are attached per unit to the pointwise model as
\texttt{cf\_*}. The 120 was chosen arbitrarily when the architecture was built, and it is a
\emph{fixed budget}, so adding features to the count model reallocates rather than accumulates.

\begin{center}\small
\begin{tabular}{lr}
\toprule
forwarded budget & micro-AUC \\
\midrule
\textbf{top-120} & \textbf{0.9632} \\
top-180 & 0.9609 \\
top-200 & 0.9618 \\
\bottomrule
\end{tabular}
\end{center}

\begin{finding}
Wider is worse. This also kills the hypothesis offered in Section~\ref{sec:negatives} for why the
slope family hurt --- that slopes crowded the budget and evicted useful features. Widening the
budget does not rescue them, so they simply do not help most units.
\end{finding}

\subsubsection*{Screening harness}

\texttt{src/screening.py} restores the convention of Section~\ref{sec:cv-harness}: every third
training CASE (not row --- striding rows would leave a case's near-duplicate hours in training,
the leak of Section~\ref{sec:state}) and 150 trees, used to rank and never to score. Calibrated
on the hourly-structure experiment it predicted $+0.0034$ where the full fit gave $+0.0025$:
right sign, right magnitude, 18s against 108s.

Stride 3 is now the standing iteration protocol for a second reason: at full stride the working
set crosses this machine's 26GB and the run begins paging. Measured, L5 at 1.7M rows takes 53s
and L7 at 5.2M rows takes 6{,}537s --- \textbf{123$\times$ for 3$\times$ the data}, a cliff rather
than a curve. Stride is held constant across arms, so deltas remain comparable while levels sit
$\sim$0.002 below full stride.
